\documentclass[preview]{standalone}
\usepackage{amsmath}
\usepackage{booktabs}
\usepackage{amssymb}
\usepackage{setspace}
\usepackage{bm}
\let\oldcdot\cdot
\usepackage{breqn}
\let\cdot\oldcdot
\begin{document}
Let $\phi(x)$ be the level-set function and
 \begin{align*}\raisebox{\depth}{$\chi$}(x) &= \begin{cases} 1 & \mbox{if}\quad \phi(x) < 0 \\ k_{air}/k_0 & \mbox{if}\quad \phi(x)  \ge 0 \end{cases},\end{align*}
where the weak-material contrast is the physically motivated ratio of the thermal conductivity of air to that of aluminium, $k_{air}/k_0 = 0.026/200 \approx 1.3 \times 10^{-4}$ (numerically $1.3 \times 10^{-4}$, the value taken by the characteristic function).
The design domain is the normalized unit square $D = [0,1] \times [0,1]$ — the cross-section of the spreader plate of a liquid-cooled power module, scaled to unit size for benchmark comparison.  The bottom edge carries the power chip: a uniform heat flux $q = 10^{3}$ W/m$^{2}$ enters through the central segment $\Gamma_q$ of the bottom edge, of length $0.4$, and the rest of the bottom edge $\Gamma_a$ is adiabatic.  The TOP edge $\Gamma_p$ is the wall of a cooling pipe (a cold plate / cold rail): it is held at the prescribed reference temperature $u_C = 0$ K, the zero of the temperature scale, i.e. an isothermal heat sink of negligible thermal resistance.  The two remaining edges form $\Gamma_c$, where the spreader exchanges heat with ambient air at $u_{amb} = 298$ K through a Robin condition with effective convective heat transfer coefficient $h = 25$ W/(m$^{2}$ K).  This is the cross-section of an electric-vehicle traction-inverter module: the dies load the baseplate, the opposite face is bolted to a water-glycol cold plate, and the sides see the air inside the housing.
The Lagrangian of the cold-plate heat-sink problem is
\begin{dmath*}
\mathcal L(\left[ u\right],\left[ \delta u\right]) = J(u) -  \int_{\Gamma_q} q \delta u \,ds_x +  \int_{D} k \raisebox{\depth}{$\chi$} \left(  \nabla u  \cdot  \nabla \delta u  \right) \,dx + \int_{\Gamma_c} h \left( u - u_{amb}\right) \delta u \,ds_x ,
\end{dmath*}
with
\begin{align*}
J(u) &=  \int_{D} \ell \raisebox{\depth}{$\chi$} \,dx + \int_{\Gamma_q} q u \,ds_x,
 \end{align*}
where The Lagrange multiplier $\ell$ weights the material cost; it is a balance parameter of the optimisation, calibrated by the pipeline (a geometric scan locating the knee of the volume-fraction curve), not part of the problem statement.
The state $\left[ u\right] \in \mathbb P_{2}(D)$ with $u = u_C$ on the cooling wall $\Gamma_p$ is the solution of the linear problem
\begin{dmath*}
-  \int_{\Gamma_q} q \delta u \,ds_x +  \int_{D} k \raisebox{\depth}{$\chi$} \left(  \nabla u  \cdot  \nabla \delta u  \right) \,dx + \int_{\Gamma_c} h u \delta u \,ds_x = \int_{\Gamma_c} h u_{amb} \delta u \,ds_x \qquad \forall\; {\delta u\in \mathbb P_{2}(D)}, \quad \delta u = 0 \ \text{on } \Gamma_p .
\end{dmath*}
The adjoint state $\left[ \tilde u\right]$ to the direct state $\left[ u\right]$ is the solution of
\begin{dmath*}
\int_{\Gamma_q} q \delta u \,ds_x +  \int_{D} k \raisebox{\depth}{$\chi$} \left(  \nabla \delta u  \cdot  \nabla \tilde u  \right) \,dx + \int_{\Gamma_c} h \delta u \tilde u \,ds_x = 0\qquad \forall\; {\delta u\in \mathbb P_{2}(D)}, \quad \delta u = 0 \ \text{on } \Gamma_p ,
\end{dmath*}
where the adjoint load $\int_{\Gamma_q} q \delta u \,ds_x$ comes from the partial derivative $\partial_{u} J(\delta u)$ of the temperature integral in the objective, and the adjoint state satisfies the homogeneous Dirichlet condition $\tilde u = 0$ on $\Gamma_p$.
In order to compute the shape derivative, we consider a shape perturbation $\mathbf x = T_t(\mathbf X) = \mathbf X + t \mathbf  V(\mathbf X)$ with a suitable velocity field $\mathbf V$, vanishing on the fixed bottom edge and on the cooling wall.  The perturbed Lagrangian $\mathcal G$ using the Lagrangian state $\left[ u\right]^{L}$ is given by
\begin{dmath*}
\mathcal G(t,\left[ u\right]^{L},\left[ \delta u\right]^{L}) = \int_{D} \ell \raisebox{\depth}{$\chi$}  \operatorname{det}\mathbf F( t )  \,dx + \int_{\Gamma_q} q u^{L} \,ds_x -  \int_{\Gamma_q} q  \delta u^{L} \,ds_x +  \int_{D} k \raisebox{\depth}{$\chi$}  \operatorname{det}\mathbf F( t )  \left( \left(  \nabla u^{L} \cdot  \mathbf F^{-1}(t) \right) \cdot \left(  \nabla \delta u^{L} \cdot  \mathbf F^{-1}(t) \right) \right) \,dx + \int_{\Gamma_c} h \left( u^{L} - u_{amb}\right) \delta u^{L} \,J_{s}(t) \,ds_x ,
\end{dmath*}
where $\mathbf F(t) = \mathbf I + t \nabla \mathbf V$ is the deformation gradient and $J_{s}(t)$ the boundary Jacobian of the moving Robin boundary $\Gamma_c$.
For the direct state $\left[ u\right]$ and the adjoint state $\left[ \tilde u\right]$, the volume expression of the shape derivative, including the curvature term on the moving Robin boundary, is given by
\begin{dmath*}
DJ(\Omega)(\mathbf V) = \int_{D} \left(\operatorname{div} {\mathbf V} \left(  \nabla u  \cdot  \nabla \tilde u  \right) - \left(  \nabla \tilde u  \cdot \left(  \nabla u  \cdot  \nabla {\mathbf V}  \right) \right) - \left(  \nabla u  \cdot \left(  \nabla \tilde u  \cdot  \nabla {\mathbf V}  \right) \right)\right) k \raisebox{\depth}{$\chi$} \,dx + \int_{D} \ell \raisebox{\depth}{$\chi$} \operatorname{div} {\mathbf V} \,dx + \int_{\Gamma_c} \kappa h \left( u - u_{amb}\right) \tilde u \left(\mathbf V \cdot \mathbf n\right) \,ds_x ,
\end{dmath*}
where $\kappa$ is the curvature of $\Gamma_c$ and $\mathbf n$ its outward unit normal.  The Robin edges $\Gamma_c$ (left and right) are design boundaries — their shape is optimised — while the bottom edge ($\Gamma_q$ together with $\Gamma_a$) and the isothermal cooling wall $\Gamma_p$ are fixed boundaries whose position is prescribed by the module layout.
The initial design is an I-BEAM (工字形) without holes, the natural skeleton of a spreader plate: a top flange of the full width $1$ and thickness $0.2$ lying along the cooling wall, a bottom flange of length $0.5$ and thickness $0.2$ under the chip footprint, and a web of width $0.2$ running the full height between them,
\begin{align*}
\phi_0(x,y) = \min\left( \left| y - 0.9 \right| - 0.1 ,\ \max\left( \left| y - 0.1 \right| - 0.1 ,\ \left| x - 0.5 \right| - 0.25 \right),\ \left| x - 0.5 \right| - 0.1 \right).
\end{align*}
The cooling wall $\Gamma_p$ is a fixed boundary: its position and the contact area with the conducting material are prescribed by the device layout, so the initial design places solid material along it (the full-width top flange) and the material in contact with it cannot detach during the level-set evolution.  The prescribed temperature $u_C$ is imposed solely through the Dirichlet value of the field declaration and appears nowhere else in the formulation.
\end{document}
